\section{Applications}\label{sec:applications}

\subsection{Gaming}\label{sec:applications_gaming}
Multiplayer video games extensively use random mechanics that directly impact player experiences and in-game economies. From critical hit calculations to loot drops and lootbox mechanics, players invest significant time and resources based on these probabilities. The fairness of these systems directly affects player satisfaction and game balance. With tiers in one library, a game can use free local randomness for cosmetic effects, the \texttt{quick} tier for frequent combat rolls and the \texttt{timelock} tier for rare drops and prizes.

This need for verifiable randomness extends naturally into gambling applications, where every slot machine pull, card shuffle, and dice roll must be demonstrably random and tamper-proof to maintain player trust. The stakes are often significant, making manipulation attempts a serious concern.

\subsection{NFTs}\label{sec:applications_nfts}
Non-Fungible Tokens (NFTs), such as those from the Bored Ape Yacht Club (BAYC)~\cite{boredapeyachtclub}, are typically generated by combining various image layers to create unique digital collectibles. Because of the Onchain Randomness Problem these are frequently created off-chain by developers and then uploaded, a process that can be manipulated. Verifiable onchain randomness enables provably fair and transparent minting.

\subsection{Proof of Stake Consensus}\label{sec:applications_defi}
Randomness is extensively utilized in Proof of Stake (PoS) consensus mechanisms. For instance, Solana employs random selection to determine the next leader in its consensus process~\cite{yakovenko2025solana}. This method ensures a fair and unpredictable rotation among validators.

\subsection{Artificial Intelligence}\label{sec:applications_security}
Random values are essential in neural networks such as large language models because they help the system learn. One of the most important uses of randomness is weight initialization, where the initial values of connections are set unpredictably so that different neurons learn different features~\cite{glorot2010understanding}.
